\documentclass[10pt,a4paper]{amsart}
\usepackage[margin=19mm]{geometry}
\usepackage{amsmath,amssymb,mathtools}
\usepackage[scr=rsfs,cal=euler]{mathalfa}
\usepackage{xfrac}
\usepackage[unicode,hidelinks,bookmarks=false]{hyperref}
\newcommand{\ie}{i.e.\ }
\newcommand{\cf}{cf.\ }
\newcommand{\resp}{respectively\ }
\newcommand{\Subsection}{\subsection}
\providecommand{\nobreakdash}{\nobreak\hbox{-}}
\setlength{\emergencystretch}{2em}
\multlinegap0pt
\usepackage{bm}
\usepackage{bbm}
\usepackage{dsfont}
\usepackage{xspace}
\usepackage{cancel}
\usepackage{mathtools}
\usepackage{amsthm}
\makeatletter
\@ifundefined{thm}{\newtheorem{thm}{Thm}}{}
\@ifundefined{lem}{\newtheorem{lem}[thm]{Lemma}}{}
\@ifundefined{prop}{\newtheorem{prop}[thm]{Proposition}}{}
\makeatother
\newcommand{\R}{\mathbb{R}}
\newcommand{\rd}{\mathrm{d}}
\title[Analysis of a one-dimensional biofilm model]{Analysis of a one-dimensional biofilm model}
\author{Patrick Guidotti, Christoph Walker}
\date{Source: arXiv:2505.23545v1 (CC BY 4.0)}
\begin{document}
\maketitle

\noindent\emph{Source and licence.} Analysis of a one-dimensional biofilm model. Version arXiv:2505.23545v1 (CC BY 4.0), distributed under the Creative Commons Attribution 4.0 International licence, as recorded in the arXiv licence metadata for this exact version and shown on the abstract page https://arxiv.org/abs/2505.23545v1. Full rights notice: licenses/arxiv-2505.23545-CC-BY-4.0.txt.\par\smallskip

\noindent\textit{Editorial scope.} Counted unit: the Prop whose source label is \texttt{P2} in the article (source0.tex lines 1317--1327, source label \texttt{P2}), delivered together with the article's own complete original proof of it (source0.tex lines 1328--1423, one \texttt{proof} environment of 96 lines). No mathematics is written, repaired or re-derived: statement and proof are the article's own text line for line. Required context retained uncounted: GeneralAssumptions (lines 254--257); P1 (lines 825--855); ccc1 (lines 843--848); BBex (lines 1049--1055); gg (lines 1214--1216); BBeex (lines 1225--1232); BBeex4 (lines 1227--1231); L10 (lines 1236--1271); c1 (lines 1242--1248); c2 (lines 1244--1247). Every definition, display and label of those lines is retained unchanged. Omitted from this package and not redistributed: 1--253, 258--824, 849--1048, 1056--1213, 1217--1224, 1232--1235, 1248--1316, 1424--1600. Nothing inside a retained excerpt is omitted or altered: every retained body is delivered exactly as the article's lines are written, so no omission receipt of any kind accompanies this package. Citations: the retained excerpts carry no citation command at all, so no bibliography entry of the article is transcribed and no reference is redistributed. Reference adaptation: none was needed and none was made; every reference inside the delivered unit resolves inside it, so no unresolved reference or citation is produced. Printed slips kept literal: no printed word, symbol, hypothesis or number of the counted statement or of its proof was altered. Layout and class: the article's own class is \texttt{amsart} with packages including times, amsmath, cancel, stmaryrd, graphicx, amsfonts, enumitem, amssymb, color, mathrsfs, hyperref, cleveref, enumerate, lmodern; the wrapper is the lane's pinned \texttt{amsart} editorial wrapper at 10pt on A4, which restates the theorem-like environments the retained text needs and adds the article's own macro definitions that the retained text needs (\texttt{\textbackslash R}, \texttt{\textbackslash rd}), with extra wrapper packages (bm, bbm, dsfont, xspace, cancel, mathtools). The delivered numbering is the unit's own. Assets: no retained span contains a figure environment, a graphics-inclusion command, a PDF-page inclusion, a table or a TikZ picture; no image file is copied into this package, the package declares no asset dependency and no image file is redistributed with it. Licence: version arXiv:2505.23545v1 of this article is distributed under the Creative Commons Attribution 4.0 International licence, as recorded in the arXiv licence metadata for this exact version and shown on the abstract page https://arxiv.org/abs/2505.23545v1. A full rights notice is delivered at licenses/arxiv-2505.23545-CC-BY-4.0.txt. Characters: every retained excerpt is pure ASCII, so the delivered engine, XeLaTeX with its default Latin Modern text font, reports no missing character.
\begin{equation}\label{GeneralAssumptions}
g\in {\rm C}^{1-}(\R), \qquad r\in
{\rm C}^{1}(\R)\ \text{ with }\ sr(s)> 0\ \text{ for }\ s\not= 0,
\end{equation}

\begin{prop}\label{P1}
Assume \eqref{GeneralAssumptions} with $r'\ge 0$. For each $h>0$ there exists a
unique solution $$u=u[h]\in {\rm C}^2 \bigl( [0,1]\bigr)$$ of the
boundary value problem
\begin{subequations}\label{BBBB}
\begin{align}
\frac{\kappa}{h^2} u_{yy}&= r(u),\quad  0<y<1,\label{BBBB1}\\
u_y(0)&=0,\quad u(1)+\frac{L\kappa}{\kappa_L h}u_y(1)=c_*.\label{BBBB2}
\end{align}
\end{subequations}
It holds that
\begin{align}\label{C}
  u(y)=c_*-\frac{Lh}{\kappa_L}\int_0^1
  r \bigl( u(\eta)\bigr)\,\rd \eta-\frac{h^2}{\kappa}
  \int_y^1\int_0^\eta r \bigl( u(\rho)\bigr)\,\rd\rho \,\rd \eta,\quad y\in [0,1].
\end{align}
Moreover, one also has that
\begin{subequations}\label{ccc}
\begin{align}
&0< u(0)\le u(y)\le u(1)< c_*,\quad  y\in [0,1],\label{ccc1}\\
&0\le u_y(y)\le u_y(1)=\frac{\kappa_L}{L\kappa}h \big(c_*-u(1)\big),\quad 
y\in [0,1],\label{ccc2}\\
&0\le u_{yy}(y)\le \frac{r(c_*)}{\kappa}h^2,\quad  y\in [0,1], \label{ccc3}
\end{align}
\end{subequations}
and
$$
\Bigl[ h\mapsto u[h]\Bigr]\in {\rm C}^1\big((0,\infty),
{\rm C}^2([0,1])\big).
$$
\end{prop}

\begin{subequations}\label{BBex}
\begin{align}
\kappa c_{zz}&= r(c),\quad 0<z<h,\label{BBex2}\\
c_z(0)&=0,\quad c(h)+\frac{L\kappa}{\kappa_L}c_z(h)=c_*, \label{BBex4}\\
&\int_0^h g\big(r(c(z))\big)\,\rd z=0. \label{BBex6}
\end{align}
 \end{subequations}

\begin{align}\label{gg}
g(s):=\alpha(s-b),\quad s\in [0,r(c_*)],
\end{align}

\begin{subequations}
\label{BBeex}
\begin{align}
\kappa c_{zz}(z)&= r\bigl( c(z)\bigr),\quad 0<z<h,\label{BBeex2}\\
c_z(0)=0,\quad c(h)&=c_*-\frac{L b}{\kappa_L}h,\quad
c_z(h)=\frac{b}{\kappa}h.\label{BBeex4} 
\end{align}
 \end{subequations}

\begin{lem}\label{L10}
Let $r\in C^1(\mathbb{R})$ be bounded, strictly increasing,
and satisfy $s r(s)\ge 0$ for $s\in \mathbb{R}$. Given any
$c_0\ge 0$ there is a unique global solution $c(\cdot,c_0)\in
{\rm C}^2 \bigl( [0,\infty)\bigr) $ to the initial value
problem
\begin{subequations}
\label{c1}
\begin{align}
\kappa\, c_{zz}&= r( c ),\quad z>0,\label{c2}\\
c(0)&=c_0,\quad  c_z(0)=0. \label{c3} 
\end{align}
 \end{subequations}
The mappings $c(\cdot,c_0)$ and $c_z(\cdot,c_0)$ are increasing and
$c(\cdot,0)\equiv 0$. In addition,
$$
\bigl[ c_0\mapsto c(\cdot,c_0)\bigr]
\in{\rm C}^1\bigl([0,\infty),{\rm C}^2([0,R])\bigr)$$
for each $R>0$. Moreover, $w:=\frac{\partial}{\partial c_0} c(\cdot,c_0)\in
{\rm C}^2\bigl([0,\infty)\bigr)$ satisfies
\begin{subequations}
\label{w1}
\begin{align}
\kappa\, w_{zz}&= r'\bigl( c(\cdot,c_0)\bigr) w(z),\quad
z>0,\label{w2}\\  w(0)&=1,\quad w_z(0)=0, \label{w3} 
\end{align}
\end{subequations}
and it holds that $w(z)\ge 1$ as well as $w_z(z)\ge 0$ for $z\ge0$.
The function $M$, defined by
$$
M(z,c_0):=\left[\frac{Lb}{\kappa_L}+c_z(z,c_0)\right]
w_z(z)-\frac{r\bigl( c(z,c_0)\bigr)-b }{\kappa}w(z),\quad z\ge 0,
$$
is strictly increasing with respect to $z\ge 0$. In particular, for $r(c_0)\le b$ it holds
that $M(z,c_0)> M(0,c_0)\ge 0$ for $z> 0$. 
\end{lem}

\begin{prop}\label{P2}
Assume \eqref{GeneralAssumptions} with $r'(s)>0$ for $s\in [0,c_*]$. It holds:
\begin{itemize}
\item[(i)] If $g\in {\rm C}^{1-}\bigl( [0,\infty)\bigr)$ is
  strictly increasing and $g \bigl( r(c_*)\bigr)\leq 0$, then the equilibrium
  problem~\eqref{BBex}  has no solution.
\item[(ii)] If $g$ is given by \eqref{gg} and $r(c_*)> b$, then the
  equilibrium problem \eqref{BBex} (or, equivalently, \eqref{BBeex}) has a unique solution $(h_e,c_e)$ 
  with $h_e>0$  and $c_e\in  {\rm C}^2([0,h_e])$.
\end{itemize}
\end{prop}

\begin{proof}
(i) If $(c,h)$ is a solution of \eqref{BBex}, then $c$ is
  non-decreasing and $c(z)< c_*$ for $z\in [0,h)$. It follows that
  $g\bigl( r\bigl(c(z)\bigr)\bigr)< 0$ for $z\in [0,h)$ thanks to
  Proposition \ref{P1}, so that
  $\int_0^h g\bigl( r\bigl(c(z)\bigr)\bigr)\,\rd z<0$ and,
  consequently, that there is no equilibrium in this case.

(ii) In order to establish that there exists a unique
  equilibrium we use a shooting argument. To this end, let
  $c(\cdot,c_0)\in {\rm C}^2([0,\infty))$ be the unique solution of
  the initial value problem~\eqref{c1} provided by Lemma \ref{L10} for
  $c_0\ge 0$. Here we may assume that $r$ satisfies the assumptions
  of Lemma~\ref{L10} as only $r\big |_{[0,c_*]}$ plays a role,
  see~\eqref{ccc1}.
  To enforce the boundary conditions \eqref{BBeex4} we define the
  functions $A$ and $B$ by
  $$
  A(z,c_0):=c(z,c_0)+\frac{L b}{\kappa_L}z-c_*, \qquad
  B(z,c_0):=c_z(z,c_0)-\frac{b}{\kappa}z,
  $$
  for $c_0,z\ge 0$. Consider $c_0\in [0,c_*)$. Using Lemma
  \ref{L10} it follows that $A(0,c_0)=c_0-c_*< 0$ and that
  $A(\cdot,c_0)$ is increasing and 
  unbounded above. Thus, there is a unique $h(c_0)> 0$ such that
  $A\bigl(h(c_0), c_0\bigr)=0$. Since 
$$
\partial_z A(h(c_0), c_0)= c_z(h(c_0),c_0)+\frac{L b}{\kappa_L}\ge
\frac{L b}{\kappa_L}>0, 
$$ 
it follows from the implicit function theorem that
$h\in{\rm C}^1\bigl([0,c_*)\bigr)$ with $h(0)=\kappa_L
c_*/(Lb)$.
Then 
  $$
  B\bigl( h(\cdot),\cdot\bigr)\in{\rm C}^1\bigl([0,c_*)\bigr).
  $$
  Since $r$ is strictly increasing and $r(0)=0$, there is a unique
  $\underline{c}\in (0,c_*)$ such that $r(\underline{c})=b$. Thus, for
  $c_0\in (\underline{c},c_*)$ we have that
  $$
  c(z,c_0)\geq c(0,c_0)=c_0>\underline{c},\quad z\geq 0,
  $$
  and hence $r\bigl( c(z,c_0)\bigr)>b$ for $z>0$ since
  $r$ is increasing. This and \eqref{c2} yield that
  $$
  \partial_z
  B(z,c_0)=c_{zz}(z,c_0)-\frac{b}{\kappa}=\frac{1}{\kappa}
  \big[r\bigl(c(z,c_0)\bigr)-b\big]>0,\quad z>0,
  $$
  while $B(0,c_0)=c_z(0,c_0)=0$. Consequently, we have that $B\bigl(
  h(c_0),c_0\bigr)>0$ for $c_0\in (\underline{c},c_*)$. On the other
  hand we have that
  $$
  B(h(0),0)=-\frac{\kappa_L c_*}{L\kappa}<0,
  $$
  and, since $B\bigl( h(\cdot),\cdot\bigr)$ depends continuously on
  its argument $c_0$, there must exist some ${\sf c}_e\in (0,\underline{c})$
  such that $B\bigl (h({\sf c}_e),{\sf c}_e\bigr)=0$. By construction, we have
  that $c(\cdot,{\sf c}_e)\in{\rm C}^2\bigl([0,h_{{\sf c}_e}]\bigr)$
  and that $(h({\sf c}_e),c\bigl( \cdot,{\sf c}_e)\bigr)$
  solves \eqref{BBeex} (where $h({\sf c}_e)>0$). In fact, we claim that
  $B\bigl( h(\cdot),\cdot\bigr)$ is strictly increasing
  on~$[0,\underline{c}]$ so that it can only possess a single zero and 
  the equilibrium is thus unique. To see this, we first differentiate
  the identity $A(h(c_0),c_0)=0$ with respect to $c_0$ and use
  Lemma~\ref{L10} to obtain 
\begin{align}\label{i1}
\left(\frac{Lb}{\kappa_L}+ c_z(h(c_0),c_0)\right)\frac{d}{d c_0}h(c_0)=-w(h(c_0))\le -1,
\end{align}
where $w=\frac{\partial}{\partial{c_0}} c(\cdot,c_0)$ as in Lemma \eqref{L10}. Next,
  we observe that
  \begin{equation*}
    \frac{d}{d c_0}  B\bigl( h(c_0), c_0\bigr)=w_z\bigl(
    h(c_0)\bigr)+\left[ c_{zz}\bigl(h(c_0),c_0\bigr)-\frac{b}{\kappa}\right]
    \frac{d}{dc_0}h(c_0),
  \end{equation*}
  and hence, using \eqref{c2} and \eqref{i1}, that
  \begin{align*}
    &\left(\frac{Lb}{\kappa_L}+c_z\bigl( h(c_0),c_0\bigr)\right)
    \frac{d}{dc_0}  B\bigl( h(c_0),c_0\bigr) \\
    &\quad =\left(\frac{Lb}{\kappa_L}+c_z\bigl(
      h(c_0),c_0\bigr)\right)
    w_z(h(c_0))-\frac{1}{\kappa}\bigl[r\bigl(c(h(c_0),c_0)\bigr)-b\bigr]
    w\bigl( h(c_0)\bigr)\\
    &\quad=M\bigl( h(c_0),c_0\bigr)>0,
  \end{align*}
  where we used that $M$, which was defined in Lemma \ref{L10}, is
  strictly positive since $r(c_0)\le b$ for $c_0\in[0,\underline{c}]$. Since Lemma
  \ref{L10}
  also ensures that $\frac{Lb}{\kappa_L}+c_z\bigl(
  h(c_0),c_0\bigr)>0$, we indeed have that $B\bigl(
  h(\cdot),\cdot\bigr)$ is strictly increasing on
  $[0,\underline{c}]$. Thus there is a single solution of
  \eqref{BBeex}, as claimed.
\end{proof}

\end{document}
